\documentclass{article}

\usepackage{amsmath, amsthm, amssymb, amsfonts}
\usepackage{thmtools}
\usepackage{setspace}
\usepackage{geometry}
\usepackage[utf8]{inputenc}
\usepackage[english]{babel}
\usepackage{framed}
\usepackage[dvipsnames]{xcolor}
\usepackage{tcolorbox}

\colorlet{LightGray}{White!90!Periwinkle}
\colorlet{LightOrange}{Orange!15}
\colorlet{LightGreen}{Green!15}
\definecolor{LightBlue}{RGB}{137, 207, 240}

\newcommand{\HRule}[1]{\rule{\linewidth}{#1}}
\newcommand{\norm}[1]{\left \lVert #1 \right \rVert}

\declaretheoremstyle[name=Theorem,]{thmsty}
\declaretheorem[style=thmsty]{theorem}
\tcolorboxenvironment{theorem}{colback=LightGray}

\declaretheoremstyle[name=Proposition,]{prosty}
\declaretheorem[style=prosty,numberlike=theorem]{proposition}
\tcolorboxenvironment{proposition}{colback=LightOrange}

\declaretheoremstyle[name=Principle,]{prcpsty}
\declaretheorem[style=prcpsty,numberlike=theorem]{principle}
\tcolorboxenvironment{principle}{colback=LightGreen}

\declaretheoremstyle[name=Definition,]{defsty}
\declaretheorem[style=defsty]{definition}
\tcolorboxenvironment{definition}{colback=LightBlue!50}

\setstretch{1.2}
\geometry{
    textheight=9in,
    textwidth=5.5in,
    top=1in,
    headheight=12pt,
    headsep=25pt,
    footskip=30pt
}

% ------------------------------------------------------------------------------

\begin{document}

% ------------------------------------------------------------------------------
% Cover Page and ToC
% ------------------------------------------------------------------------------

\title{ \normalsize \textsc{}
\\ [2.0cm]
\HRule{1.5pt} \\
\LARGE \textbf{\uppercase{Math 291}
\HRule{2.0pt} \\ [0.6cm] \LARGE{Honors Multivariable Calculus} \vspace*{10\baselineskip}}
}
\date{}
\author{\textbf{Jeffrey Lu}}


% ------------------------------------------------------------------------------

\section*{Chapter 2 — Description of Motion}

\subsection*{2.1 Functions from $\mathbb{R}$ to $\mathbb{R}^n$}

\begin{align*}
    \mathbf{x}(t)      & = (x_1(t), x_2(t), \ldots, x_n(t))                             \\
    \mathbf{x}'(t)     & = \frac{d\mathbf{x}}{dt} = (x_1'(t), x_2'(t), \ldots, x_n'(t)) \\
    \|\mathbf{x}'(t)\| & = \sqrt{(x_1'(t))^2 + (x_2'(t))^2 + \cdots + (x_n'(t))^2}      \\
    s(t)               & = \int_{t_0}^{t} \|\mathbf{x}'(u)\| \, du
\end{align*}

\subsection*{2.1.5 Arc Length Parameterization}

\begin{align*}
    \frac{d\mathbf{x}}{ds} & = \mathbf{T}(s)          \\
    \frac{d\mathbf{T}}{ds} & = \kappa(s)\mathbf{N}(s)
\end{align*}

\subsection*{2.1.6 Frenet–Serret Formulas}

\begin{align*}
    \mathbf{T}'(s) & = \kappa(s)\mathbf{N}(s)                         \\
    \mathbf{N}'(s) & = -\kappa(s)\mathbf{T}(s) + \tau(s)\mathbf{B}(s) \\
    \mathbf{B}'(s) & = -\tau(s)\mathbf{N}(s)
\end{align*}



\subsection*{2.1.7 Curvature and Torsion}

\begin{align*}
    \kappa & = \frac{\|\mathbf{v} \times \mathbf{a}\|}{\|\mathbf{v}\|^3}                          \\
    \tau   & = \frac{\mathbf{a}' \cdot (\mathbf{v} \times \mathbf{a})}{\|\mathbf{v}\|^6 \kappa^2}
    = -\frac{\mathbf{x}''' \cdot (\mathbf{x}'' \times \mathbf{x}')}{\|\mathbf{x}'\|^6 \kappa^2}
\end{align*}


The Darboux vector $\boldsymbol{\omega}$ describes the instantaneous angular velocity
of the moving Frenet frame $\{\mathbf{T}, \mathbf{N}, \mathbf{B}\}$ along a space curve.

\[
    \frac{d\mathbf{T}}{ds} = \boldsymbol{\omega} \times \mathbf{T}, \quad
    \frac{d\mathbf{N}}{ds} = \boldsymbol{\omega} \times \mathbf{N}, \quad
    \frac{d\mathbf{B}}{ds} = \boldsymbol{\omega} \times \mathbf{B}
\]

where $\boldsymbol{\omega}$ is defined by
\[
    \boldsymbol{\omega} = \tau \mathbf{T} + \kappa \mathbf{B}
\]

Hence, substituting into the Frenet–Serret formulas:

\begin{align}
    \frac{d\mathbf{T}}{ds} & = \kappa \mathbf{N},                    \\
    \frac{d\mathbf{N}}{ds} & = -\kappa \mathbf{T} + \tau \mathbf{B}, \\
    \frac{d\mathbf{B}}{ds} & = -\tau \mathbf{N}.
\end{align}

\subsection*{2.1.8 Helix Parameterization}

If curvature $\kappa$ and torsion $\tau$ are constants, then
\begin{align*}
    \mathbf{x}(t) & = \mathbf{x}(0)
    + \frac{\kappa}{\kappa^2 + \tau^2}
    \begin{bmatrix}
        (1 - \cos(\sqrt{\kappa^2 + \tau^2}\,t))\,\mathbf{u}_1
        - \sin(\sqrt{\kappa^2 + \tau^2}\,t)\,\mathbf{u}_2
    \end{bmatrix}
    + t\,\frac{\tau}{\sqrt{\kappa^2 + \tau^2}}\,\mathbf{u}_3
\end{align*}

Constants for radius and pitch:
\begin{align*}
    r = \frac{\kappa}{\kappa^2 + \tau^2}, \quad
    c = \sqrt{\kappa^2 + \tau^2}, \quad
    b = \frac{\kappa}{\kappa^2 + \tau^2}
\end{align*}

$$\left| \frac{1-n}{n} + 1 \right| = \left| \frac{1-n}{n} + \frac{n}{n} \right| = \left| \frac{1 - n + n}{n} \right| = \left| \frac{1}{n} \right| $$Since $n$ is a positive integer, $|1/n|$ is just $1/n$. **This is the key\!** We found that the distance from our sequence to the limit is just **$1/n$**. Let's check this: * At $n=10$, the term is $\frac{1-10}{10} = -0.9$. The distance to -1 is $|-0.9 - (-1)| = 0.1$. And $1/n$ is $1/10 = 0.1$. It works\! * At $n=100$, the term is $\frac{1-100}{100} = -0.99$. The distance to -1 is $|-0.99 - (-1)| = 0.01$. And $1/n$ is $1/100 = 0.01$. It works again\!  2\. Win the Game (Finding $N$) Now we can respond to the skeptic's challenge. * **Skeptic's Challenge:** "Make your distance less than my tolerance $\epsilon$\!" * **Our Task:** We need to find $n$ large enough so that:$$
$$
\end{document}
